(1) 无限深圆阱内部满足 $-\hbar^2\nabla^2\psi/(2m)=E\psi$。取 $\psi=R(r)e^{i\ell\theta}$，则 $r^2R''+rR'+(k^2r^2-\ell^2)R=0$，$k=\sqrt{2mE}/\hbar$。原点的平方可积正则解取 $J_{|\ell|}(kr)$，排除在原点发散的Neumann支。

(2) 由硬壁 $R(a)=0$，有 $ka=j_{|\ell|,n}$，其中 $j_{|\ell|,n}$ 是第 $n$ 个正Bessel零点。能级 $\boxed{E_{\ell n}=\hbar^2j_{|\ell|,n}^2/(2ma^2)}$。$\ell=0$ 非角向简并，$\pm\ell\ne0$ 两支等能。基态为 $\ell=0,n=1$，是无节点轴对称模。

(3) 利用 $\int_0^arJ_{|\ell|}^2(j_{|\ell|,n}r/a)dr=a^2J_{|\ell|+1}^2(j_{|\ell|,n})/2$，归一复波函数为 $\psi_{\ell n}=J_{|\ell|}(j_{|\ell|,n}r/a)e^{i\ell\theta}/[a\sqrt\pi|J_{|\ell|+1}(j_{|\ell|,n})|]$。径向概率密度 $p_r(r)=2rJ_{|\ell|}^2(j_{|\ell|,n}r/a)/[a^2J_{|\ell|+1}^2(j_{|\ell|,n})]$，积分至 $a$ 为1。

(4) $E_{\ell n}\propto a^{-2}$，所以壁面向外的广义力 $\boxed{F_a=-\partial_aE=2E/a>0}$。扩大势阱减少量子动能，因此粒子对边界施加向外压力；按圆周长平均的二维边缘压力大小为 $F_a/(2\pi a)=E/(\pi a^2)$。这与二维量子维里缩放一致。
